% Auto-generated from raw.txt
% Usage:
%   \vocabentry{word}{/ipa/ (pos)}{definition}{example}{note}

\vocabentry{Aesthetic}
{/esˈθetɪk/ (adj)}
{Concerned with beauty or the appreciation of beauty.}
{The building was designed with both practical and aesthetic considerations in mind.}
{Đây là một tính từ dùng để chỉ các yếu tố thuộc về thẩm mỹ hoặc cảm thụ vẻ đẹp.}

\vocabentry{Artifact}
{/ˈɑːrtɪfækt/ (n)}
{An object made by a human being.}
{These gold artifacts were recovered from a shipwreck.}
{Trong quyền động vật, có thể chỉ các công cụ săn bắn cổ.}

\vocabentry{Heritage}
{/ˈherɪtɪdʒ/ (n)}
{Valued objects and qualities such as cultural traditions and historic buildings passed down from previous generations.}
{We must work together to preserve our national heritage for future generations.}
{Đây là một danh từ dùng để chỉ di sản văn hóa hoặc truyền thống.}

\vocabentry{Masterpiece}
{/ˈmæstərpiːs/ (n)}
{A work of outstanding artistry, skill, or workmanship.}
{Leonardo da Vinci's "Mona Lisa" is considered a masterpiece of the Renaissance.}
{Đây là một danh từ dùng để chỉ một kiệt tác hoặc tác phẩm nghệ thuật để đời.}

\vocabentry{Curation}
{/kjʊˈreɪʃn/ (n)}
{The action or process of selecting, organizing, and looking after the items in a collection or exhibition.}
{The success of the museum depends on the careful curation of its permanent collection.}
{Đây là một danh từ dùng để chỉ công tác giám tuyển, tổ chức triển lãm.}

\vocabentry{Abstract}
{/ˈæbstrækt/ (adj)}
{Existing in thought or as an idea but not having a physical or concrete existence.}
{Abstract art often uses shapes and colors to represent emotions rather than real objects.}
{Đây là một tính từ dùng để chỉ nghệ thuật trừu tượng.}

\vocabentry{Renaissance}
{/ˌrenəˈsɑːns/ (n)}
{The revival of art and literature under the influence of classical models in the 14th–16th centuries.}
{The Renaissance was a period of great intellectual and artistic achievement in Europe.}
{Đây là một danh từ dùng để chỉ thời kỳ Phục hưng.}

\vocabentry{Contemporary}
{/kənˈtempəreri/ (adj)}
{Living or occurring at the same time; belonging to or occurring in the present.}
{The gallery features works by both established and contemporary artists.}
{Đây là một tính từ dùng để chỉ tính đương đại, hiện đại.}

\vocabentry{Folklore}
{/ˈfoʊklɔːr/ (n)}
{The traditional beliefs, customs, and stories of a community, passed through the generations by word of mouth.}
{Much of the author's work is based on Irish folklore and mythology.}
{Đây là một danh từ dùng để chỉ văn hóa dân gian hoặc truyền thuyết.}

\vocabentry{Avant-garde}
{/ˌævɑ̃ː ˈɡɑːrd/ (adj)}
{New and unusual or experimental ideas, especially in the arts, or the people introducing them.}
{The filmmaker is known for his avant-garde style and unconventional storytelling.}
{Đây là một tính từ dùng để chỉ phong cách tiên phong hoặc phá cách.}

\vocabentry{Calligraphy}
{/kəˈlɪɡrəfi/ (n)}
{Decorative handwriting or handwritten lettering.}
{Traditional Vietnamese calligraphy is often practiced during the Lunar New Year.}
{Đây là một danh từ dùng để chỉ nghệ thuật thư pháp.}

\vocabentry{Iconography}
{/ˌaɪkəˈnɑːɡrəfi/ (n)}
{The visual images and symbols used in a work of art or the study or interpretation of these.}
{The religious iconography in the medieval cathedral is incredibly detailed.}
{Đây là một danh từ dùng để chỉ hệ biểu tượng hoặc sự nghiên cứu hình tượng nghệ thuật.}

\vocabentry{Metaphor}
{/ˈmetəfɔːr/ (n)}
{A figure of speech in which a word or phrase is applied to an object or action.}
{The journey in the story is a metaphor for the character's personal growth.}
{Phép ẩn dụ trong các câu chuyện lịch sử.}

\vocabentry{Sculpture}
{/ˈskʌlptʃər/ (n)}
{The art of making two- or three-dimensional representative or abstract forms, especially by carving stone or wood.}
{The city park is home to several impressive modern sculptures.}
{Đây là một danh từ dùng để chỉ nghệ thuật điêu khắc.}

\vocabentry{Canvas}
{/ˈkænvəs/ (n)}
{A strong, coarse unbleached cloth used as a surface for oil painting.}
{The painter spent weeks preparing the large canvas for his next work.}
{Đây là một danh từ dùng để chỉ khung vải vẽ hoặc bề mặt hội họa.}

\vocabentry{Genre}
{/ˈʒɑːnrə/ (n)}
{A category of artistic composition characterized by similarities in form, style, or subject matter.}
{Landscape painting has been a popular genre for centuries.}
{Đây là một danh từ dùng để chỉ thể loại nghệ thuật.}

\vocabentry{Patronage}
{/ˈpætrənɪdʒ/ (n)}
{The support given by a patron (a wealthy or influential person) to an artist or organization.}
{Without the patronage of the wealthy Medici family, many Renaissance works would not exist.}
{Đây là một danh từ dùng để chỉ sự bảo trợ nghệ thuật.}

\vocabentry{Acoustics}
{/əˈkuːstɪks/ (n)}
{The properties or qualities of a room or building that determine how sound is transmitted in it.}
{The concert hall is famous for its perfect acoustics.}
{Đây là một danh từ dùng để chỉ âm học hoặc độ vang của âm thanh.}

\vocabentry{Portrait}
{/ˈpɔːrtreɪt/ (n)}
{A painting, drawing, photograph, or engraving of a person.}
{The royal family commissioned an official portrait to celebrate the anniversary.}
{Đây là một danh từ dùng để chỉ tranh chân dung.}

\vocabentry{Etching}
{/ˈetʃɪŋ/ (n)}
{An engraved plate or print prepared with the use of acid.}
{The gallery has a fine collection of 18th-century etchings.}
{Đây là một danh từ dùng để chỉ kỹ thuật khắc axit hoặc bản khắc hội họa.}

\vocabentry{Performance}
{/pərˈfɔːrməns/ (n)}
{An act of staging or presenting a play, concert, or other form of entertainment.}
{The actors gave a powerful performance that moved the audience to tears.}
{Đây là một danh từ dùng để chỉ buổi biểu diễn nghệ thuật.}

\vocabentry{Harmony}
{/ˈhɑːrməni/ (n)}
{The combination of simultaneously sounded musical notes to produce chords.}
{The choir sang in perfect harmony throughout the entire concert.}
{Đây là một danh từ dùng để chỉ sự hòa hợp hoặc hòa âm trong âm nhạc.}

\vocabentry{Rhythm}
{/ˈrɪðəm/ (n)}
{A strong, regular, repeated pattern of movement or sound.}
{The fast rhythm of the traditional drums made everyone want to dance.}
{Đây là một danh từ dùng để chỉ nhịp điệu.}

\vocabentry{Exhibition}
{/ˌeksɪˈbɪʃn/ (n)}
{A public display of works of art or items of interest, held in an art gallery or museum.}
{I'm going to see an exhibition of Japanese prints this weekend.}
{Đây là một danh từ dùng để chỉ triển lãm.}

\vocabentry{Composition}
{/ˌkɑːmpəˈzɪʃn/ (n)}
{The way in which a whole or mixture is made up; a work of music, literature, or art.}
{The painter carefully considered the composition of the scene before starting to paint.}
{Đây là một danh từ dùng để chỉ bố cục tác phẩm hoặc bản soạn nhạc.}

\vocabentry{Impressionism}
{/ɪmˈpreʃənɪzəm/ (n)}
{A style of painting that seeks to capture the visual impression of the moment.}
{Claude Monet is one of the most famous figures associated with Impressionism.}
{Đây là một danh từ dùng để chỉ trường phái Ấn tượng.}

\vocabentry{Symbolism}
{/ˈsɪmbəlɪzəm/ (n)}
{The use of symbols to represent ideas or qualities.}
{The white dove in the painting is rich with symbolism, representing peace.}
{Đây là một danh từ dùng để chỉ chủ nghĩa biểu tượng hoặc hệ thống biểu tượng.}

\vocabentry{Perspective}
{/pərˈspektɪv/ (n)}
{The art of drawing solid objects on a two-dimensional surface so as to give the right impression of height, width, and depth.}
{Renaissance artists revolutionized painting by using mathematical perspective.}
{Đây là một danh từ dùng để chỉ luật phối cảnh trong hội họa.}

\vocabentry{Narrative}
{/ˈnærətɪv/ (n)}
{A spoken or written account of connected events; a story.}
{The author's narrative style is engaging and easy to follow.}
{Lời kể/Văn tự sự lịch sử.}

\vocabentry{Architecture}
{/ˈɑːrkɪtektʃər/ (n)}
{The art or practice of designing and constructing buildings.}
{The city is famous for its blend of modern and traditional architecture.}
{Đây là một danh từ dùng để chỉ ngành kiến trúc.}

\vocabentry{Sublime}
{/səˈblaɪm/ (adj)}
{Of such excellence, grandeur, or beauty as to inspire great admiration and awe.}
{The view of the Alps at sunset was truly sublime.}
{Đây là một tính từ dùng để chỉ vẻ đẹp cao cả, tráng lệ.}

\vocabentry{Ensemble}
{/ɑ̃ːnˈsɑːmbl/ (n)}
{A group of musicians, actors, or dancers who perform together.}
{The chamber music ensemble performed works by Bach and Mozart.}
{Đây là một danh từ dùng để chỉ một dàn nhạc, đoàn diễn hoặc nhóm nghệ sĩ.}

\vocabentry{Visual}
{/ˈvɪʒuəl/ (adj)}
{Relating to seeing or sight.}
{Visual arts include painting, sculpture, and photography.}
{Đây là một tính từ dùng để chỉ các yếu tố thuộc về thị giác.}

\vocabentry{Literary}
{/ˈlɪtəreri/ (adj)}
{Concerning the writing, study, or content of literature.}
{The magazine published a fascinating interview with a famous literary figure.}
{}

\vocabentry{Authenticity}
{/ˌɔːθenˈtɪsəti/ (n)}
{The quality of being authentic or genuine.}
{Experts were called in to verify the authenticity of the ancient scroll.}
{Đây là một danh từ dùng để chỉ tính nguyên bản hoặc sự xác thực của tác phẩm.}

\vocabentry{Mosaic}
{/moʊˈzeɪɪk/ (n)}
{A picture or pattern produced by arranging together small colored pieces of hard material, such as stone, tile, or glass.}
{The walls of the ancient villa were decorated with beautiful mosaics.}
{Đây là một danh từ dùng để chỉ nghệ thuật khảm đá hoặc tranh ghép mảnh.}

\vocabentry{Tapestry}
{/ˈtæpɪstri/ (n)}
{A piece of thick textile fabric with pictures or designs formed by weaving colored weft threads.}
{The historic castle is filled with ornate tapestries depicting heroic battles.}
{Đây là một danh từ dùng để chỉ thảm thêu hoặc sự phong phú đa dạng.}

\vocabentry{Melody}
{/ˈmelədi/ (n)}
{A sequence of musical notes that form the main tune of a song or piece of music.}
{The melody is simple, but it stays in your head for hours.}
{Đây là một danh từ dùng để chỉ giai điệu.}

\vocabentry{Anthropology}
{/ˌænθrəˈpɑːlədʒi/ (n)}
{The study of human societies and cultures and their development.}
{Anthropology helps us understand the diverse cultural practices around the world.}
{Đây là một danh từ dùng để chỉ nhân chủng học, nghiên cứu văn hóa.}

\vocabentry{Conservation}
{/ˌkɑːnsərˈveɪʃn/ (n)}
{The protection of plants and animals and natural areas from the damaging effects of human activity.}
{Conservation experts are working to clean the ancient statues.}
{Đây là một danh từ dùng để chỉ sự bảo tồn động vật hoang dã và môi trường sống của chúng.}

\vocabentry{Diversity}
{/daɪˈvɜːrsəti/ (n)}
{The state of being diverse; variety; especially cultural diversity.}
{Cultural diversity is a source of strength and creativity for our society.}
{Đây là một danh từ dùng để chỉ sự đa dạng, đặc biệt là đa dạng văn hóa.}

\vocabentry{Appropriation}
{/əˌproʊpriˈeɪʃn/ (n)}
{The action of taking something for one's own use, typically without the owner's permission.}
{The artist was accused of cultural appropriation for using traditional indigenous motifs.}
{Trong nghệ thuật, đây là sự chiếm dụng văn hóa hoặc vay mượn hình ảnh.}

\vocabentry{Fusion}
{/ˈfjuːʒn/ (n)}
{The process or result of joining two or more things together to form a single entity.}
{Jazz fusion combines elements of jazz, rock, and funk.}
{Đây là một danh từ dùng để chỉ sự giao thoa hoặc kết hợp các phong cách nghệ thuật/ẩm thực.}

\vocabentry{Installation}
{/ˌɪnstəˈleɪʃn/ (n)}
{An artistic genre of three-dimensional works that are often site-specific and designed to transform the perception of a space.}
{The latest installation at the museum uses lights and sounds to create an immersive forest.}
{Đây là một danh từ dùng để chỉ nghệ thuật sắp đặt.}

\vocabentry{Motif}
{/moʊˈtiːf/ (n)}
{A decorative design or pattern; a distinctive feature or dominant idea in an artistic or literary composition.}
{The recurring flower motif in her paintings represents growth and renewal.}
{Đây là một danh từ dùng để chỉ họa tiết hoặc chủ đề chủ đạo.}

\vocabentry{Revival}
{/rɪˈvaɪvl/ (n)}
{An improvement in the condition or strength of something; an instance of something becoming popular or active again.}
{There has been a significant revival of interest in traditional folk music.}
{Đây là một danh từ dùng để chỉ sự phục hưng hoặc sự sống lại của một phong cách.}

\vocabentry{Symmetry}
{/ˈsɪmətri/ (n)}
{The quality of being made up of exactly similar parts facing each other or around an axis.}
{The classical temple is a perfect example of balance and symmetry.}
{Đây là một danh từ dùng để chỉ sự đối xứng.}

\vocabentry{Universal}
{/ˌjuːnɪˈvɜːrsl/ (adj)}
{Relating to or done by all people or things in the world or in a particular group.}
{Music is often called the universal language of mankind.}
{Đây là một tính từ dùng để chỉ tính phổ quát toàn cầu.}

\vocabentry{Vernacular}
{/vərˈnækjələr/ (adj)}
{(Of language or architecture) spoken or used by the ordinary people in a particular country or region.}
{The museum focuses on the vernacular architecture of the local farming community.}
{Đây là một tính từ dùng để chỉ tính bản địa, dân dã.}

\vocabentry{Vibrant}
{/ˈvaɪbrənt/ (adj)}
{Full of energy and enthusiasm; (of color) bright and striking.}
{The city has a vibrant arts scene with many small galleries and theaters.}
{Đây là một tính từ dùng để chỉ sự sôi nổi, rực rỡ.}

\vocabentry{Abstraction}
{/æbˈstrækʃn/ (n)}
{The quality of dealing with ideas rather than events.}
{Her work is characterized by a high degree of abstraction and minimalism.}
{Đây là một danh từ dùng để chỉ sự trừu tượng hóa.}

\vocabentry{Bas-relief}
{/ˌbɑː rɪˈliːf/ (n)}
{A type of sculpture that has less depth to the faces and forms than they actually have.}
{The temple walls are decorated with bas-reliefs depicting scenes from ancient myths.}
{Đây là một danh từ dùng để chỉ nghệ thuật chạm nổi hoặc phù điêu.}

\vocabentry{Catalogue}
{/ˈkætəlɔːɡ/ (n)}
{A complete list of items, typically one in alphabetical or other systematic order.}
{I bought a catalogue of the exhibition to learn more about the artists.}
{Trong nghệ thuật là danh mục triển lãm.}

\vocabentry{Ceramics}
{/səˈræmɪks/ (n)}
{Pots and other articles made from clay hardened by heat.}
{The museum has a world-class collection of Chinese ceramics.}
{Đây là một danh từ dùng để chỉ đồ gốm sứ.}

\vocabentry{Chorus}
{/ˈkɔːrəs/ (n)}
{A large organized group of singers, especially one that performs together with an orchestra or opera company.}
{The finale of the opera featured a magnificent performance by the chorus.}
{Đây là một danh từ dùng để chỉ dàn hợp xướng.}

\vocabentry{Cinematography}
{/ˌsɪnəməˈtɑːɡrəfi/ (n)}
{The art of making motion pictures.}
{The movie won an award for its stunning cinematography and use of natural light.}
{Đây là một danh từ dùng để chỉ nghệ thuật điện ảnh hoặc kỹ thuật quay phim.}

\vocabentry{Classic}
{/ˈklæsɪk/ (adj)}
{Judged over a period of time to be of the highest quality and outstanding of its kind.}
{"Pride and Prejudice" is a classic example of the English novel.}
{Đây là một tính từ dùng để chỉ tính kinh điển.}

\vocabentry{Collage}
{/kəˈlɑːʒ/ (n)}
{A piece of art made by sticking various different materials such as photographs and pieces of paper or fabric on to a backing.}
{Students created a collage using old magazines and newspapers.}
{Đây là một danh từ dùng để chỉ nghệ thuật cắt dán tranh.}

\vocabentry{Commission}
{/kəˈmɪʃn/ (v)}
{Formally chosen to produce a work of art or to perform a particular duty.}
{The city council commissioned a local artist to create a mural for the library.}
{Đây là một động từ dùng để chỉ việc đặt làm, đặt hàng một tác phẩm nghệ thuật.}

\vocabentry{Cosmopolitan}
{/ˌkɑːzməˈpɑːlɪtən/ (adj)}
{Including people from many different countries and cultures.}
{London is a cosmopolitan city where you can experience a wide range of cultures.}
{Đây là một tính từ dùng để chỉ tính đa văn hóa, mang tầm thế giới.}

\vocabentry{Critic}
{/ˈkrɪtɪk/ (n)}
{A person who expresses an unfavorable opinion of something.}
{The art critic gave the new exhibition a very positive review.}
{Trong nghệ thuật là nhà phê bình.}

\vocabentry{Criticism}
{/ˈkrɪtɪsɪzəm/ (n)}
{The expression of disapproval of someone or something based on perceived faults or mistakes.}
{Literary criticism helps us understand the deeper meanings of a text.}
{Đây là một danh từ dùng để chỉ sự phê bình hoặc lý luận phê bình.}

\vocabentry{Culture}
{/ˈkʌltʃər/ (n)}
{The arts and other manifestations of human intellectual achievement regarded collectively.}
{Reading books and going to museums are great ways to engage with different cultures.}
{Đây là một danh từ dùng để chỉ văn hóa.}

\vocabentry{Curator}
{/kjʊˈreɪtər/ (n)}
{A keeper or custodian of a museum or other collection.}
{The curator organized a special tour of the museum's hidden treasures.}
{Đây là một danh từ dùng để chỉ người giám tuyển hoặc quản thủ bảo tàng.}

\vocabentry{Custom}
{/ˈkʌstəm/ (n)}
{A traditional and widely accepted way of behaving or doing something that is specific to a particular society, place, or time.}
{It is a local custom to exchange gifts during the festival.}
{Đây là một danh từ dùng để chỉ phong tục, tập quán.}

\vocabentry{Decor}
{/deɪˈkɔːr/ (n)}
{The furnishing and decoration of a room.}
{The restaurant's decor is elegant and modern.}
{Đây là một danh từ dùng để chỉ phong cách trang trí nội thất.}

\vocabentry{Design}
{/dɪˈzaɪn/ (n/v)}
{A plan or drawing produced to show the look and function or workings of a building, garment, or other object before it is built or made.}
{The innovative design of the chair makes it both comfortable and stylish.}
{Đây là một danh từ/động từ dùng để chỉ bản thiết kế hoặc việc thiết kế.}

\vocabentry{Dialect}
{/ˈdaɪəlekt/ (n)}
{A particular form of a language which is peculiar to a specific region or social group.}
{The characters in the play speak in a rural dialect that can be difficult to understand.}
{Đây là một danh từ dùng để chỉ tiếng địa phương hoặc phương ngữ.}

\vocabentry{Dimension}
{/daɪˈmenʃn/ (n)}
{An aspect or feature of a situation, problem, or thing.}
{The large-scale sculpture adds a new dimension to the public square.}
{Trong nghệ thuật dùng để chỉ chiều không gian.}

\vocabentry{Drama}
{/ˈdrɑːmə/ (n)}
{A play for theater, radio, or television.}
{She decided to study drama at university to become a professional actress.}
{Đây là một danh từ dùng để chỉ kịch nghệ hoặc sự kịch tính.}

\vocabentry{Drawing}
{/ˈdrɔːɪŋ/ (n)}
{A picture or diagram made with a pencil, pen, or crayon rather than paint.}
{The architect produced several detailed drawings of the proposed building.}
{Đây là một danh từ dùng để chỉ tranh vẽ nét hoặc ký họa.}

\vocabentry{Echo}
{/ˈekoʊ/ (n/v)}
{A sound or series of sounds caused by the reflection of sound waves from a surface back to the listener.}
{The themes of the poem echo those found in his earlier work.}
{Trong nghệ thuật dùng để chỉ sự vang vọng hoặc mô phỏng lại.}

\vocabentry{Elite}
{/eˈliːt/ (n/adj)}
{A select group that is superior in terms of ability or qualities to the rest of a group or society.}
{Opera was once considered an entertainment only for the social elite.}
{Trong văn hóa dùng để chỉ tầng lớp tinh hoa.}

\vocabentry{Eloquent}
{/ˈeləkwənt/ (adj)}
{Fluent or persuasive in speaking or writing.}
{The speaker gave an eloquent defense of the importance of art education.}
{Đây là một tính từ dùng để chỉ sự hùng hồn, có sức thuyết phục.}

\vocabentry{Empathy}
{/ˈempəθi/ (n)}
{The ability to understand and share the feelings of another.}
{Literature can help develop empathy by allowing us to see the world through others' eyes.}
{Nghệ thuật thường khơi gợi sự đồng cảm.}

\vocabentry{Engraving}
{/ɪnˈɡreɪvɪŋ/ (n)}
{A print made from an engraved plate, block, or other surface.}
{The book is illustrated with several fine wood engravings.}
{Đây là một danh từ dùng để chỉ kỹ thuật điêu khắc khắc gỗ/kim loại hoặc bản in khắc.}

\vocabentry{Epic}
{/ˈepɪk/ (adj/n)}
{A long poem, typically one derived from ancient oral tradition, narrating the deeds and adventures of heroic or legendary figures.}
{The "Iliad" is one of the most famous epic poems in history.}
{Đây là một tính từ/danh từ dùng để chỉ tính sử thi hoặc bản anh hùng ca.}

\vocabentry{Ethics}
{/ˈeθɪks/ (n)}
{Moral principles that govern a person's behavior or the conducting of an activity.}
{The film explores the complex ethics of genetic engineering.}
{Văn hóa và nghệ thuật thường liên quan đến đạo đức.}

\vocabentry{Ethnocentric}
{/ˌeθnoʊˈsentrɪk/ (adj)}
{Evaluating other peoples and cultures according to the standards of one's own culture.}
{We must avoid making ethnocentric judgments when studying other societies.}
{Đây là một tính từ dùng để chỉ tính duy chủng tộc, thiên kiến văn hóa.}

\vocabentry{Evolution}
{/ˌevəˈluːʃn/ (n)}
{The gradual development of something, especially from a simple to a more complex form.}
{The museum documents the evolution of painting styles over the last five centuries.}
{Dùng để chỉ sự tiến hóa của nghệ thuật/văn hóa.}

\vocabentry{Exotic}
{/ɪɡˈzɑːtɪk/ (adj)}
{Originating in or characteristic of a distant foreign country.}
{The botanic garden features a wide variety of exotic plants from around the world.}
{Đây là một tính từ dùng để chỉ vẻ đẹp kỳ lạ, ngoại lai.}

\vocabentry{Experimental}
{/ɪkˌsperɪˈmentl/ (adj)}
{(Of a new invention or product) based on untested ideas or techniques and not yet established or finalized.}
{The experimental music performance used many unusual sounds and instruments.}
{Đây là một tính từ dùng để chỉ tính thực nghiệm hoặc phá cách trong nghệ thuật.}

\vocabentry{Expression}
{/ɪkˈspreʃn/ (n)}
{The action of making known one's thoughts or feelings.}
{Artistic expression is an essential part of the human experience.}
{Đây là một danh từ dùng để chỉ sự biểu đạt.}

\vocabentry{Expressionism}
{/ɪkˈspreʃənɪzəm/ (n)}
{A style of painting, music, or drama in which the artist or writer seeks to express emotional experience.}
{Munch's "The Scream" is a classic example of Expressionism.}
{Đây là một danh từ dùng để chỉ trường phái Biểu hiện.}

\vocabentry{Fable}
{/ˈfeɪbl/ (n)}
{A short story, typically with animals as characters, conveying a moral.}
{Aesop's fables are still widely read and loved today.}
{Đây là một danh từ dùng để chỉ truyện ngụ ngôn.}

\vocabentry{Fiction}
{/ˈfɪkʃn/ (n)}
{Literature in the form of prose, especially short stories and novels, that describes imaginary events and people.}
{Science fiction often explores futuristic technology and its impact on society.}
{Đây là một danh từ dùng để chỉ văn học hư cấu.}

\vocabentry{Figurative}
{/ˈfɪɡjərətɪv/ (adj)}
{Representing forms that are recognizably derived from real life.}
{Figurative art focuses on representing the human body and other real objects.}
{Đây là một tính từ dùng để chỉ nghệ thuật hình thể, ngược lại với trừu tượng.}

\vocabentry{Finale}
{/fɪˈnæli/ (n)}
{The last part of a piece of music, an entertainment, or a public event, especially when particularly dramatic or exciting.}
{The fireworks display provided a spectacular finale to the festival.}
{Đây là một danh từ dùng để chỉ phần kết, đoạn cuối.}

\vocabentry{Fine arts}
{/ˌfaɪn ˈɑːrts/ (n)}
{Creative art, especially visual art whose products are to be appreciated primarily or solely for their imaginative, aesthetic, or intellectual content.}
{She decided to study fine arts at the university to specialize in painting.}
{Đây là một cụm danh từ dùng để chỉ mỹ thuật.}

\vocabentry{Form}
{/fɔːrm/ (n)}
{The visible shape or configuration of something.}
{The sculptor is interested in the human form and its movements.}
{Trong nghệ thuật dùng để chỉ hình thức hoặc trường phái.}

\vocabentry{Formal}
{/ˈfɔːrml/ (adj)}
{Done in accordance with rules of convention or etiquette.}
{The formal garden is designed with strict symmetry and order.}
{Dùng để chỉ phong cách trang trọng hoặc thuộc về hình thức.}

\vocabentry{Fresco}
{/ˈfreskoʊ/ (n)}
{A painting done rapidly in watercolor on wet plaster on a wall or ceiling.}
{Michelangelo's frescoes on the ceiling of the Sistine Chapel are world-famous.}
{Đây là một danh từ dùng để chỉ tranh bích họa vẽ trên tường khi vôi còn ướt.}

\vocabentry{Gala}
{/ˈɡeɪlə/ (n)}
{A social occasion with special entertainments or performances.}
{The charity held a gala dinner to raise money for the local arts center.}
{Đây là một danh từ dùng để chỉ buổi dạ tiệc nghệ thuật.}

\vocabentry{Gallery}
{/ˈɡæləri/ (n)}
{A room or building for the display or sale of works of art.}
{The modern art gallery has several rooms dedicated to contemporary sculpture.}
{Đây là một danh từ dùng để chỉ phòng trưng bày hoặc triển lãm tranh.}

\vocabentry{Geographic}
{/ˌdʒiːəˈɡræfɪk/ (adj)}
{Relating to geography.}
{Geographic isolation has helped to preserve the unique traditions of the island.}
{Văn hóa thường gắn liền với địa lý.}

\vocabentry{Global}
{/ˈɡloʊbl/ (adj)}
{Relating to the whole world; worldwide.}
{In today's global culture, we are increasingly influenced by ideas from all over the world.}
{Trong văn hóa là tính toàn cầu.}

\vocabentry{Gothic}
{/ˈɡɑːθɪk/ (adj)}
{Of or in the style of architecture prevalent in western Europe in the 12th–16th centuries.}
{The city's cathedral is a magnificent example of Gothic architecture.}
{Đây là một tính từ dùng để chỉ phong cách kiến trúc hoặc văn học Gô-tích.}

\vocabentry{Gourmet}
{/ˈɡʊrmeɪ/ (adj)}
{Relating to good food and drink.}
{The festival features a wide variety of gourmet food from local producers.}
{Ẩm thực là một phần của văn hóa.}

\vocabentry{Grandeur}
{/ˈɡrændʒər/ (n)}
{Splendor and impressiveness, especially of appearance or style.}
{We were awed by the grandeur of the ancient palace ruins.}
{Đây là một danh từ dùng để chỉ sự tráng lệ, vĩ đại.}

\vocabentry{Graphic}
{/ˈɡræfɪk/ (adj)}
{Relating to visual art, especially drawing, design, and lettering.}
{Graphic design is a vital part of modern advertising and media.}
{Đây là một tính từ dùng để chỉ đồ họa.}

\vocabentry{Guiding}
{/ˈɡaɪdɪŋ/ (adj)}
{Showing the way.}
{Respect for elders is one of the guiding principles of their culture.}
{Thường dùng cho "guiding principles" - các nguyên tắc văn hóa chủ đạo.}

\vocabentry{Habitat}
{/ˈæbɪtæt/ (n)}
{The natural home or environment of an animal, plant, or other organism.}
{Protecting the natural habitat is essential for preserving the culture of indigenous peoples.}
{Môi trường sống cũng ảnh hưởng đến văn hóa bản địa.}

\vocabentry{Harmonious}
{/hɑːrˈmoʊniəs/ (adj)}
{Tuneful; not discordant; or (of a whole) forming a pleasing or consistent whole.}
{The colors in the painting are balanced and harmonious.}
{Đây là một tính từ dùng để chỉ sự hài hòa.}

\vocabentry{Haze}
{/heɪz/ (n)}
{A slight obscuration of the lower atmosphere.}
{The artist used a soft haze to create a dreamy and atmospheric landscape.}
{Trong nghệ thuật dùng để mô tả không khí mờ ảo.}

\vocabentry{Heirloom}
{/ˈerluːm/ (n)}
{A valuable object that has belonged to a family for several generations.}
{This antique watch is a family heirloom that has been passed down for over a hundred years.}
{Đây là một danh từ dùng để chỉ vật gia bảo.}

\vocabentry{Hero}
{/ˈhɪroʊ/ (n)}
{A person who is admired or idealized for courage, outstanding achievements, or noble qualities.}
{The epic poem tells the story of a legendary hero who saved his people.}
{Hình tượng nhân vật trung tâm trong văn hóa.}

\vocabentry{Historic}
{/hɪˈstɔːrɪk/ (adj)}
{Famous or important in history.}
{The historic town center is a popular destination for tourists.}
{Đây là một tính từ dùng để chỉ tính lịch sử hoặc địa danh lịch sử.}

\vocabentry{Historical}
{/hɪˈstɔːrɪkl/ (adj)}
{Concerning past events.}
{The novel is based on meticulous historical research.}
{Đây là một tính từ dùng để chỉ những gì thuộc về lịch sử.}

\vocabentry{Humanism}
{/ˈhjuːmənɪzəm/ (n)}
{A rationalistic outlook or system of thought attaching prime importance to human rather than divine or supernatural matters.}
{Humanism was a central theme of Renaissance art and literature.}
{Đây là một danh từ dùng để chỉ chủ nghĩa nhân văn.}

\vocabentry{Icon}
{/ˈaɪkɑːn/ (n)}
{A person or thing regarded as a representative symbol or as worthy of veneration.}
{Marilyn Monroe remains a global fashion and beauty icon.}
{Đây là một danh từ dùng để chỉ biểu tượng văn hóa.}

\vocabentry{Identity}
{/aɪˈdentəti/ (n)}
{The fact of being who or what a person or thing is.}
{Maintaining your traditional customs is important for preserving your cultural identity.}
{Văn hóa là cốt lõi của bản sắc.}

\vocabentry{Ideology}
{/ˌaɪdiˈɑːlə dʒi/ (n)}
{A system of ideas and ideals, especially one that forms the basis of economic or political theory and policy.}
{The prevailing ideology of the time is often reflected in its popular literature.}
{Hệ tư tưởng ảnh hưởng đến nghệ thuật.}

\vocabentry{Illuminate}
{/ɪˈluːmɪneɪt/ (v)}
{Light up; or help to clarify or explain.}
{The medieval manuscript is beautifully illuminated with gold leaf and bright colors.}
{Trong hội họa cổ là nghệ thuật trang trí bản thảo.}

\vocabentry{Illusion}
{/ɪˈluːʒn/ (n)}
{A thing that is or is likely to be wrongly perceived or interpreted by the senses.}
{The artist used clever lighting to create the illusion of depth on a flat surface.}
{Nghệ thuật thường tạo ra sự ảo giác thị giác.}

\vocabentry{Illustration}
{/ˌɪləˈstreɪʃn/ (n)}
{A picture illustrating a book, newspaper, etc.}
{The children's book features several colorful and engaging illustrations.}
{Đây là một danh từ dùng để chỉ tranh minh họa.}

\vocabentry{Imagery}
{/ˈɪmɪdʒəri/ (n)}
{Visually descriptive or figurative language, especially in a literary work.}
{The poet uses vivid imagery to describe the beauty of the autumn landscape.}
{Đây là một danh từ dùng để chỉ ngôn ngữ hình ảnh hoặc hệ thống hình ảnh.}

\vocabentry{Immersive}
{/ɪˈmɜːrsɪv/ (adj)}
{(Of a computer display or system) generating a three-dimensional image which appears to surround the user.}
{The virtual reality exhibition provided a truly immersive experience of ancient Rome.}
{Nghệ thuật trải nghiệm nhập vai.}

\vocabentry{Indigenous}
{/ɪnˈdɪdʒənəs/ (adj)}
{Originating or occurring naturally in a particular place; native.}
{The museum has a collection of art and artifacts from the indigenous people of Australia.}
{Văn hóa bản địa.}

\vocabentry{Individualism}
{/ˌɪndɪˈvɪdʒuəlɪzəm/ (n)}
{A social theory favoring freedom of action for individuals over collective or state control.}
{The rise of individualism led to a greater focus on personal expression in art.}
{Chủ nghĩa cá nhân trong sáng tạo nghệ thuật.}

\vocabentry{Infusion}
{/ɪnˈfjuːʒn/ (n)}
{The introduction of a new element or quality into something.}
{The infusion of new ideas from abroad helped to revitalize the local arts scene.}
{Sự truyền cảm hứng hoặc pha trộn văn hóa.}

\vocabentry{Inherent}
{/ɪnˈherənt/ (adj)}
{Existing in something as a permanent, essential, or characteristic attribute.}
{A love of music seems to be inherent in almost every human culture.}
{Cố hữu, vốn có.}

\vocabentry{Initial}
{/ɪˈnɪʃl/ (adj)}
{Existing or occurring at the beginning.}
{The artist's initial sketches show how the idea for the painting developed.}
{Các bản phác thảo ban đầu.}

\vocabentry{Initiative}
{/ɪˈnɪʃətɪv/ (n)}
{The ability to assess and initiate things independently.}
{The government launched a new initiative to support young artists and musicians.}
{Các sáng kiến văn hóa.}

\vocabentry{Innovation}
{/ˌɪnəˈveɪʃn/ (n)}
{The action or process of innovating.}
{Technological innovation has led to the creation of many new forms of digital art.}
{Sự đổi mới nghệ thuật.}

\vocabentry{Insight}
{/ˈaɪsaɪt/ (n)}
{A deep understanding of a person or thing.}
{The novel provides a fascinating insight into life in 19th-century Russia.}
{Sự thấu hiểu thông qua nghệ thuật.}

\vocabentry{Inspirational}
{/ˌɪnspəˈreɪʃənl/ (adj)}
{Providing or showing creative or spiritual inspiration.}
{The teacher's speech was truly inspirational and encouraged many students to pursue art.}
{Đây là một tính từ dùng để chỉ sự truyền cảm hứng.}

\vocabentry{Instrument}
{/ˈɪnstrəmənt/ (n)}
{A tool or implement, especially one for precision work; or a musical instrument.}
{She plays several musical instruments, including the piano and the violin.}
{Nhạc cụ hoặc công cụ sáng tác.}

\vocabentry{Integral}
{/ˈɪntɪɡrəl/ (adj)}
{Necessary to make a whole complete; essential or fundamental.}
{Traditional music is an integral part of the local festival celebrations.}
{Thiết yếu, không thể tách rời.}

\vocabentry{Integration}
{/ˌɪntɪˈɡreɪʃn/ (n)}
{The action or process of integrating.}
{The successful integration of immigrants into the local community is a slow process.}
{Sự hòa nhập văn hóa.}

\vocabentry{Integrity}
{/ɪnˈteɡrəti/ (n)}
{The quality of being honest and having strong moral principles.}
{The artist is known for his artistic integrity and refusal to compromise his vision.}
{Sự chính trực trong nghệ thuật.}

\vocabentry{Intellect}
{/ˈɪntəlekt/ (n)}
{The faculty of reasoning and understanding objectively.}
{Engaging with challenging art and literature can help to stimulate the intellect.}
{Trí tuệ.}

\vocabentry{Interaction}
{/ˌɪntərˈækʃn/ (n)}
{Reciprocal action or influence.}
{The interaction between different cultures can lead to great creativity and innovation.}
{Sự tương tác văn hóa.}

\vocabentry{Intercultural}
{/ˌɪntərˈkʌltʃərəl/ (adj)}
{Relating to or involving different cultures.}
{The university offers an intercultural studies program to promote mutual understanding.}
{Giao lưu văn hóa.}

\vocabentry{Interpretation}
{/ɪnˌtɜːrprɪˈteɪʃn/ (n)}
{The action of explaining the meaning of something.}
{Every viewer will have their own interpretation of the abstract painting.}
{Sự diễn giải hoặc cảm thụ nghệ thuật.}

\vocabentry{Intuition}
{/ˌɪntuˈɪʃn/ (n)}
{The ability to understand something immediately, without the need for conscious reasoning.}
{The artist relies on her intuition when choosing which colors to use.}
{Trực giác sáng tạo.}

\vocabentry{Invoice}
{/ˈɪnvɔɪs/ (n)}
{A list of goods sent or services provided, with a statement of the sum due for these.}
{The gallery sent an invoice to the collector for the two paintings he bought.}
{Trong kinh doanh nghệ thuật.}

\vocabentry{Irony}
{/ˈaɪrəni/ (n)}
{The expression of one's meaning by using language that normally signifies the opposite.}
{There is a sense of tragic irony in the play's ending.}
{Sự mỉa mai, trớ trêu trong nghệ thuật/văn học.}

\vocabentry{Irreplaceable}
{/ˌɪrɪˈpleɪsəbl/ (adj)}
{Too valuable or rare to be replaced.}
{The ancient manuscripts are irreplaceable treasures of our cultural history.}
{Không thể thay thế, cực kỳ quý giá.}

\vocabentry{Journal}
{/ˈdʒɜːrnl/ (n)}
{A newspaper or magazine that deals with a particular subject or professional activity.}
{She publishes her research in a prestigious academic journal of art history.}
{Tạp chí chuyên ngành nghệ thuật.}

\vocabentry{Justice}
{/ˈdʒʌstɪs/ (n)}
{Just behavior or treatment.}
{Many modern plays explore themes of social justice and equality.}
{Công lý là chủ đề phổ biến trong nghệ thuật.}

\vocabentry{Kitsch}
{/kɪtʃ/ (n)}
{Art, objects, or design considered to be in poor taste because of excessive garishness or sentimentality.}
{Some people love the kitsch decor of the retro cafe, while others find it tacky.}
{Đây là một danh từ dùng để chỉ loại nghệ thuật "sến", lòe loẹt nhưng đôi khi có chủ ý.}

\vocabentry{Knowledge}
{/ˈnɑːlɪdʒ/ (n)}
{Facts, information, and skills acquired by a person through experience or education.}
{A deep knowledge of history is essential for understanding classical art.}
{Tri thức văn hóa.}

\vocabentry{Landscape}
{/ˈlændskeɪp/ (n)}
{All the visible features of an area of countryside or land.}
{The artist is famous for his beautiful landscapes of the English countryside.}
{Tranh phong cảnh.}

\vocabentry{Language}
{/ˈlæŋɡwɪdʒ/ (n)}
{The method of human communication.}
{Learning a foreign language can help you understand a different culture more deeply.}
{Ngôn ngữ là một phần của văn hóa.}

\vocabentry{Legend}
{/ˈledʒənd/ (n)}
{A traditional story sometimes popularly regarded as historical but unauthenticated.}
{The ancient ruins are the subject of many local legends and myths.}
{Truyền thuyết.}

\vocabentry{Legitimate}
{/ləˈdʒɪtɪmət/ (adj)}
{Conforming to the law or to rules.}
{The museum only acquires artifacts through legitimate and ethical channels.}
{Hợp pháp hoặc chính đáng.}

\vocabentry{Leisure}
{/ˈliːʒər/ (n)}
{Use of free time for enjoyment.}
{Going to the cinema and visiting art galleries are popular leisure activities.}
{Hoạt động thư giãn gắn với văn hóa.}

\vocabentry{Liberal}
{/ˈlɪbərəl/ (adj)}
{Open to new behavior or opinions and willing to discard traditional values.}
{The city is known for its liberal and tolerant social atmosphere.}
{Tư tưởng tự do.}

\vocabentry{Library}
{/ˈlaɪbreri/ (n)}
{A building or room containing collections of books and other materials.}
{I spent the whole afternoon studying in the university library.}
{Thư viện - nơi lưu giữ văn hóa.}

\vocabentry{Limitation}
{/ˌlɪmɪˈteɪʃn/ (n)}
{A limiting rule or circumstance; a restriction.}
{Working within certain limitations can actually help to stimulate artistic creativity.}
{Sự hạn chế sáng tạo.}

\vocabentry{Linear}
{/ˈlɪniər/ (adj)}
{Arranged in or extending along a straight or nearly straight line.}
{The drawing uses strong linear elements to create a sense of movement.}
{Trong nghệ thuật là đường nét hoặc cốt truyện tuyến tính.}

\vocabentry{Linguistic}
{/lɪŋˈɡwɪstɪk/ (adj)}
{Relating to language or linguistics.}
{The country's linguistic diversity is reflected in its many different regional dialects.}
{Thuộc về ngôn ngữ học.}

\vocabentry{Literacy}
{/ˈlɪtərəsi/ (n)}
{The ability to read and write.}
{Promoting literacy is essential for ensuring equal access to information and culture.}
{Trình độ học vấn/văn hóa.}

\vocabentry{Literature}
{/ˈlɪtrətʃər/ (n)}
{Written works, especially those considered of superior or lasting artistic merit.}
{Classical literature can provide valuable insights into the beliefs of past societies.}
{Văn học.}

\vocabentry{Logic}
{/ˈlɑːdʒɪk/ (n)}
{Reasoning conducted or assessed according to strict principles of validity.}
{There is a clear and consistent logic to the way the building is designed.}
{Logic trong bố cục nghệ thuật.}

\vocabentry{Longevity}
{/lɔːnˈdʒevəti/ (n)}
{Long life.}
{The longevity of traditional crafts depends on passing skills to younger generations.}
{Sự trường tồn của văn hóa/nghệ thuật.}

\vocabentry{Lyrical}
{/ˈlɪrɪkl/ (adj)}
{(Of literature, art, or music) expressing the writer's emotions in an imaginative and beautiful way.}
{The author's prose is so lyrical that it often feels like reading poetry.}
{Tính chất trữ tình.}

\vocabentry{Mainstream}
{/ˈmeɪnstriːm/ (adj)}
{Representing widespread current thought.}
{Independent films often provide a welcome alternative to mainstream cinema.}
{Văn hóa/nghệ thuật đại chúng.}

\vocabentry{Manual}
{/ˈmænjuəl/ (adj)}
{Relating to or done with the hands.}
{Weaving is a traditional manual skill that requires a great deal of patience.}
{Lao động/thủ công mỹ nghệ.}

\vocabentry{Manuscript}
{/ˈmænjuskrɪpt/ (n)}
{A book, document, or piece of music written by hand rather than typed or printed.}
{The library holds a rare collection of ancient medieval manuscripts.}
{Bản thảo viết tay.}

\vocabentry{Marginalized}
{/ˈmɑːrdʒɪnəlaɪzd/ (adj)}
{(Of a person, group, or concept) treated as insignificant or peripheral.}
{The exhibition aims to give a voice to marginalized communities and their stories.}
{Các nhóm bị gạt ra ngoài lề văn hóa.}

\vocabentry{Masquerade}
{/ˌmæskəˈreɪd/ (n)}
{A false show or pretense; also a masked ball.}
{The festival features a spectacular masquerade parade with colorful costumes.}
{Lễ hội hóa trang hoặc sự giả tạo.}

\vocabentry{Mechanism}
{/ˈmekənɪzəm/ (n)}
{A system of parts working together in a machine; or a process by which something takes place.}
{Museums have complex mechanisms for preserving and protecting their artifacts.}
{Cơ chế hoạt động của các thiết chế văn hóa.}

\vocabentry{Mediate}
{/ˈmiːdieɪt/ (v)}
{Intervene between people in a dispute in order to bring about an agreement.}
{Art can help to mediate cultural differences and promote mutual understanding.}
{Nghệ thuật là cầu nối trung gian.}

\vocabentry{Medieval}
{/ˌmediˈiːvl/ (adj)}
{Relating to the Middle Ages.}
{The city is famous for its well-preserved medieval architecture.}
{Thuộc về thời Trung cổ.}

\vocabentry{Medium}
{/ˈmiːdiəm/ (n)}
{The material or form used by an artist, composer, or writer.}
{Oil on canvas has been a popular medium for portrait painters for centuries.}
{Chất liệu hoặc phương tiện nghệ thuật.}

\vocabentry{Memoir}
{/ˈmemwɑːr/ (n)}
{A historical account or biography written from personal knowledge.}
{The former leader wrote a fascinating memoir about his time in office.}
{Hồi ký - một thể loại văn học.}

\vocabentry{Memory}
{/ˈmeməri/ (n)}
{The faculty by which the mind stores and remembers information.}
{Many of her paintings are based on childhood memories of her family home.}
{Nghệ thuật lưu giữ ký ức.}

\vocabentry{Metropolis}
{/məˈtrɑːpəlɪs/ (n)}
{The capital or chief city of a country or region.}
{Tokyo is a vast metropolis that offers a wide range of cultural attractions.}
{Trung tâm văn hóa đô thị.}

\vocabentry{Minimalism}
{/ˈmɪnɪməlɪzəm/ (n)}
{A trend in sculpture and painting that arose in the 1950s and used simple geometric forms.}
{Minimalism emphasizes the importance of space and simplicity in design.}
{Chủ nghĩa tối giản.}

\vocabentry{Modernization}
{/ˌmɑːrdənəˈzeɪʃn/ (n)}
{The process of adapting something to modern needs or habits.}
{The city is struggling to balance modernization with the preservation of its heritage.}
{Hiện đại hóa văn hóa.}

\vocabentry{Monarchy}
{/ˈmɑːnərki/ (n)}
{A form of government with a monarch at the head.}
{The grandeur of the palace reflects the power and wealth of the monarchy.}
{Chế độ quân chủ ảnh hưởng đến nghệ thuật cung đình.}

\vocabentry{Monolith}
{/ˈmɑːnəlɪθ/ (n)}
{A large single upright block of stone, especially one shaped into or serving as a pillar or monument.}
{The ancient temple features several massive stone monoliths.}
{Khối đá đơn độc hoặc công trình khổng lồ.}

\vocabentry{Monument}
{/ˈmɑːnjumənt/ (n)}
{A statue, building, or other structure, erected to commemorate a notable person or event.}
{The war monument was built to honor those who lost their lives in the conflict.}
{Đài tưởng niệm hoặc công trình vĩ đại.}

\vocabentry{Mural}
{/ˈmjʊrəl/ (n)}
{A painting or other work of art executed directly on a wall.}
{The neighborhood is famous for its colorful murals depicting local life.}
{Tranh tường.}

\vocabentry{Museum}
{/mjuˈziːəm/ (n)}
{A building in which objects of historical, scientific, artistic, or cultural interest are stored and exhibited.}
{The British Museum has one of the world's largest collections of ancient artifacts.}
{Bảo tàng.}

\vocabentry{Myth}
{/mɪθ/ (n)}
{A traditional story, especially one concerning the early history of a people.}
{Ancient Greek myths often feature powerful gods and heroic figures.}
{Thần thoại.}

\vocabentry{Native}
{/ˈneɪtɪv/ (adj)}
{Associated with the country, region, or circumstances of a person's birth.}
{He is a native speaker of several indigenous languages.}
{Bản địa.}

\vocabentry{Neutral}
{/ˈnuːtrəl/ (adj)}
{Not supporting or helping either side in a conflict, disagreement, etc.}
{The artist aims to present a neutral and objective view of the social issues.}
{Tính trung lập của nghệ thuật.}

\vocabentry{Nomadic}
{/noʊˈmædɪk/ (adj)}
{Living the life of a nomad; wandering.}
{The exhibition documents the traditional customs of nomadic people in Central Asia.}
{Văn hóa du mục.}

\vocabentry{Norm}
{/nɔːrm/ (n)}
{Something that is usual, typical, or standard.}
{Respect for family is a fundamental social norm in many Asian cultures.}
{Chuẩn mực văn hóa.}

\vocabentry{Novel}
{/ˈnɑːvl/ (n/adj)}
{A long story, typically about imaginary characters and events; or something new and unusual.}
{Her latest novel has been praised for its original plot and complex characters.}
{Tiểu thuyết hoặc tính mới lạ.}

\vocabentry{Nuance}
{/ˈnuːɑːns/ (n)}
{A subtle difference in or shade of meaning, expression, or sound.}
{A good actor can convey a wide range of emotions through subtle nuances in performance.}
{Sắc thái tinh tế trong nghệ thuật.}

\vocabentry{Nurture}
{/ˈnɜːrtʃər/ (v)}
{Care for and encourage the growth or development of.}
{The arts center aims to nurture young talent and encourage creativity.}
{Đây là một động từ dùng để chỉ sự nuôi dưỡng hoặc bồi dưỡng tài năng.}

\vocabentry{Objective}
{/əbˈdʒektɪv/ (adj)}
{(Of a person or their judgment) not influenced by personal feelings or opinions.}
{It is difficult to remain objective when criticizing a work of art that you love.}
{Tính khách quan.}

\vocabentry{Obsolete}
{/ˌɑːbsəˈliːt/ (adj)}
{No longer produced or used; out of date.}
{Many traditional crafts have become obsolete due to industrialization.}
{Lỗi thời, không còn dùng.}

\vocabentry{Opulence}
{/ˈɑːpjələns/ (n)}
{Great wealth or luxuriousness.}
{The opulence of the royal court is reflected in the ornate furniture and jewelry.}
{Sự sang trọng, phong phú của nghệ thuật quý tộc.}

\vocabentry{Orchestra}
{/ˈɔːrkɪstrə/ (n)}
{A large group of musicians playing various instruments together.}
{The symphony orchestra performed works by Beethoven and Brahms.}
{Dàn nhạc giao hưởng.}

\vocabentry{Ornate}
{/ɔːrˈneɪt/ (adj)}
{Made in an intricate shape or decorated with complex patterns.}
{The ancient temple features beautifully ornate carvings and mosaics.}
{Được trang trí công phu.}

\vocabentry{Outlet}
{/ˈaʊtlet/ (n)}
{A means of expressing one's talents, energy, or emotions.}
{Painting provides her with a creative outlet for her emotions.}
{Cửa ngõ/phương tiện bộc lộ sáng tạo.}

\vocabentry{Paradigm}
{/ˈpærədaɪm/ (n)}
{A typical example or pattern of something; a model.}
{The rise of digital technology has created a new paradigm for artistic production.}
{Mô hình hoặc chuẩn mực văn hóa.}

\vocabentry{Participation}
{/pɑːrˌtɪsɪˈpeɪʃn/ (n)}
{The action of taking part in something.}
{Encouraging community participation is essential for the success of the arts festival.}
{Sự tham gia hoạt động văn hóa.}

\vocabentry{Passage}
{/ˈpæsɪdʒ/ (n)}
{A short part of a book, speech, or piece of music.}
{The speaker read a beautiful passage from a poem by Robert Frost.}
{Một đoạn trích tác phẩm.}

\vocabentry{Pastoral}
{/ˈpæstərəl/ (adj)}
{(Of a work of art) portraying or evoking country life, typically in a romanticized or idealized way.}
{The artist is known for his idyllic pastoral scenes of grazing sheep and rolling hills.}
{Phong cách điền viên, đồng quê.}

